% 天王星（综述）
% license CCBYSA3
% type Wiki

本文根据 CC-BY-SA 协议转载翻译自维基百科\href{https://en.wikipedia.org/wiki/Uranus}{相关文章}。
\begin{figure}[ht]
\centering
\includegraphics[width=8cm]{./figures/fb24dd1c06777ad6.png}
\caption{由旅行者 2 号拍摄的天王星真彩图像\(^\text{[a]}\)。其苍白而柔和的外观源于覆盖在云层之上的一层薄雾。} \label{fig_Uranus_1}
\end{figure}
天王星是距离太阳的第七颗行星。它是一颗呈青色的气态冰巨行星。该行星的大部分由处于超临界相状态的水、氨和甲烷构成，天文学上将这些物质称为“冰”或挥发物。该行星的大气具有复杂的分层云结构，并在所有太阳系行星中具有最低的最低温度（49 K（−224 °C; −371 °F））。它具有 82.23° 的显著自转轴倾角，自转周期为 17 小时 14 分钟，且为逆行自转。这意味着在其围绕太阳完成一个为期 84 个地球年的公转周期期间，其两极会经历约 42 年的连续日照，随后是约 42 年的连续黑暗。

在太阳系行星中，天王星直径排名第三、质量排名第四。根据当前模型，在其挥发性地幔层内部存在一个岩质核心，外围包覆着厚厚的氢和氦大气。在其高层大气中检测到痕量的碳氢化合物（被认为通过水解作用产生）以及一氧化碳和二氧化碳（被认为源自彗星）。天王星大气中存在许多尚未解释的气候现象，例如其峰值风速 900 km/h（560 mph）\(^\text{[24]}\)、极冠的变化以及其不规则的云层结构。该行星的内部热量相较其他巨行星异常低，其原因尚不清楚。

与其他巨行星一样，天王星拥有环系统、磁层以及众多天然卫星。其极为暗淡的环系统仅反射约 2\% 的入射光。天王星的 29 颗天然卫星中包含 19 颗已知的规则卫星，其中 14 颗为小型内侧卫星。在更外侧的是该行星五颗较大的主要卫星：米兰达（Miranda）、艾丽尔（Ariel）、安布里尔（Umbriel）、泰坦妮亚（Titania）和奥伯龙（Oberon）。在距离天王星更远的轨道上则有十颗已知的不规则卫星。该行星的磁层高度不对称，且含有大量带电粒子，这可能是其环和卫星变暗的原因。

天王星肉眼可见，但亮度极低，并且相对于背景恒星移动非常缓慢，因此直到 1781 年被威廉·赫歇尔首次观测到之前，它从未被归类为一颗行星。在其发现约七十年后，人们达成共识，将其命名为希腊原始神祇之一乌拉诺斯（Uranus，Ouranos）。截至 2025 年，它仅被探测器访问过一次，即 1986 年旅行者二号（Voyager 2）飞掠该行星\(^\text{[25]}\)。尽管如今可通过望远镜分辨并观测其细节，但科研界强烈希望再次造访该行星——这一点体现于《行星科学十年调查》将拟议的“天王星轨道器与探测器任务”列为 2023–2032 年的最高优先级任务之一，以及中国国家航天局提出利用天问四号的子探测器飞掠天王星的计划\(^\text{[26]}\)。
\subsection{历史}
1781 年 3 月 13 日（天王星被发现之日），天王星的位置（以叉号标示）与古典行星一样，天王星肉眼可见，但由于其亮度极低且公转缓慢，古代观测者从未将其识别为一颗行星\(^\text{[27]}\)。威廉·赫歇尔于 1781 年 3 月 13 日首次观测到天王星，由此发现了这颗行星，这也是人类历史上首次扩展太阳系已知边界，并使天王星成为第一颗借助望远镜而被确认的行星。天王星的发现实际上使当时已知太阳系的尺度扩大了一倍，因为天王星距离太阳的平均距离约为土星的两倍。
\subsubsection{发现}
在天王星被确认是一颗行星之前，它曾被多次观测到，但一般被误认为是一颗恒星。显然，喜帕恰斯在公元前 128 年观测过它，当时他为自己的星表测量恒星的位置，该星表后来被纳入托勒密的《天文大成》中。该星表给出了处于室女座的四颗恒星的位置，它们构成一个四边形。其中有一颗恒星实际上并不存在，而在公元前 128 年 4 月，天王星正位于该位置\(^\text{[28]}\)。最早的确定无误的观测记录是在 1690 年，当时约翰·弗拉姆斯蒂德至少六次观测到它，并将其编入星表为金牛座 34 号星。詹姆斯·布拉德雷在 1748 年、1750 年和 1753 年三次观测到它。托比亚斯·迈耶于 1756 年观测到它一次。法国天文学家皮埃尔·夏尔·勒莫尼耶在 1750 年至 1769 年期间至少十二次观测到天王星，其中包括连续四个夜晚的观测\(^\text{[29]}\)。

威廉·赫歇尔于 1781 年 3 月 13 日，在英国萨默塞特郡巴斯新国王街 19 号自家花园（现为赫歇尔天文学博物馆）观测到了天王星\(^\text{[30]}\)，并于 1781 年 4 月 26 日最初将其报告为一颗彗星\(^\text{[31]}\)。使用一架自制的 6.2 英寸反射望远镜，赫歇尔“从事了一系列关于恒星视差的观测”\(^\text{[32][33]}\)。

赫歇尔在日记中写道：“在靠近金牛座 ζ 的四分仪区域……要么是一颗星云状恒星，要么可能是一颗彗星”\(^\text{[34]}\)。3 月 17 日他记道：“我寻找那颗彗星或星云状恒星，并发现它是一颗彗星，因为它改变了位置”\(^\text{[35]}\)。当他向英国皇家学会呈报他的发现时，他依然坚持认为自己发现的是一颗彗星，但也含蓄地将其与一颗行星作比较\(^\text{[32]}\)：

“当我第一次看到这颗彗星时所使用的放大倍率是 227。根据经验我知道，恒星的视直径在更高倍率下不会按比例放大，而行星则会；因此我随即将倍率调至 460 和 932，发现彗星的直径随倍率按比例增大，这正符合其并非恒星的假设，而与之比较的那些恒星的直径并未按相同比例增大。此外，这颗彗星在远超其亮度所能承受的高倍率下，显得模糊且轮廓不清，而恒星则保持了我从成千上万次观测中所熟知的那种光泽与清晰度。随后的一切证明了我的猜测是正确的，这正是我们最近观测到的那颗彗星。”

赫歇尔将他的发现通知了皇家天文学家内维尔·马斯克林，并在 1781 年 4 月 23 日收到他困惑的回信：“我不知道该如何称呼它。它既有可能是一颗沿着近圆形轨道绕太阳运行的常规行星，也有可能是一颗沿着极端椭圆轨道运行的彗星。我尚未看到它有任何彗发或彗尾”\(^\text{[36]}\)。

尽管赫歇尔继续将他的新天体描述为一颗彗星，但其他天文学家已经开始怀疑并非如此。芬兰—瑞典籍、在俄罗斯工作的天文学家安德斯·约翰·勒克塞尔是第一位计算出这一新天体轨道的人\(^\text{[37]}\)。其近乎圆形的轨道暗示它是一颗行星而非一颗彗星。柏林天文学家约翰·埃勒特·博德将赫歇尔的发现描述为“一颗移动的恒星，可被视为迄今未知的、在土星轨道之外运行的类行星天体”\(^\text{[38]}\)。博德得出结论认为，其近圆形轨道与行星轨道更为相似，而非彗星轨道\(^\text{[39]}\)。

该天体很快被接受为一颗新行星。到 1783 年，赫歇尔向皇家学会会长约瑟夫·班克斯承认了这一点：“据欧洲最杰出天文学家的观测结果显示，那颗我在 1781 年 3 月有幸向他们指出的新星，是我们太阳系中的一颗主行星”\(^\text{[40]}\)。为表彰他的成就，乔治三世国王授予赫歇尔每年 200 英镑的津贴（相当于 2023 年的 30,000 英镑）\(^\text{[41]}\)，条件是他搬到温莎，以便王室成员能够通过他的望远镜观测天体\(^\text{[42]}\)。
\begin{figure}[ht]
\centering
\includegraphics[width=8cm]{./figures/221d0cedb3ff1f58.png}
\caption{} \label{fig_Uranus_3}
\end{figure}
\subsubsection{姓名}
\begin{figure}[ht]
\centering
\includegraphics[width=6cm]{./figures/57b83b23d4c37bca.png}
\caption{威廉·赫歇尔，天王星的发现者} \label{fig_Uranus_2}
\end{figure}
“天王星”这一名称源自古希腊的天空神乌拉诺斯（古希腊语：Οὐρανός），在罗马神话中称为 Caelus，为克洛诺斯（即土星）的父亲、宙斯（即木星）的祖父以及阿瑞斯（即火星）的曾祖父，其名称在拉丁语中被写作 Uranus（IPA: [ˈuːranʊs]）\(^\text{[2]}\)。它是八大行星中唯一一个英文名称来自希腊神话人物的行星。天文学家所偏好的 “Uranus” 发音是 /ˈjʊərənəs/（YOOR-ə-nəs）\(^\text{[1]}\)，其发音方式包含英语的长 “u”，且重音落在第一音节，与拉丁语 Uranus 相同；另一种常见发音 /jʊˈreɪnəs/（yoo-RAY-nəs）则将重音置于第二音节并使用长 a，两者皆被视为可接受\(^\text{[g]}\)。

关于行星名称的共识直到发现后的近 70 年后才最终达成。在最初的讨论中，马斯基林曾请求赫歇尔“请施惠于天文学界，为你亲自发现、令我们深表感激的这颗行星赐名”\(^\text{[44]}\)。作为回应，赫歇尔决定将其命名为 Georgium Sidus（乔治之星），或称“乔治行星”，以纪念他的新庇护人乔治三世国王\(^\text{[45]}\)。他在写给约瑟夫·班克斯的信中解释了这一决定\(^\text{[40]}\)：

在古代神话时代，水星、金星、火星、木星与土星被命名为行星，是以其主要英雄与神祇之名相称。在当今更具哲学精神的时代，再以同样方式为我们新的天体命名为朱诺、帕拉斯、阿波罗或密涅瓦，则未免不太合宜。任何特殊事件或显著事物，其首要考虑应为年代记载：若在未来时代，人们询问这颗最新发现的行星是何时被发现的？那么，“在乔治三世统治时期”将是一个令人十分满意的答案。

赫歇尔提出的名称在不列颠与汉诺威之外并不受欢迎，于是很快出现了替代方案。天文学家热罗姆·拉朗德建议将其命名为 Herschel 以纪念发现者\(^\text{[46]}\)。瑞典天文学家 Erik Prosperin 提出了 Astraea、Cybele（后来成为小行星名称）以及 Neptune（后来成为下一颗被发现的行星的名称）。来自哥廷根的格奥尔格·李希滕贝格也支持 Astraea（作 Austräa），但该女神传统上与处女座而非金牛座相关。其他天文学家支持 Neptune，希望以此纪念英国皇家海军在美国独立战争期间的胜利，他们提出将新行星命名为 Neptune George III 或 Neptune Great Britain；这一折衷方案也由莱克塞尔所建议\(^\text{[37][47]}\)。丹尼尔·伯努利则提出了 Hypercronius 与 Transaturnis。Minerva 亦曾被提议为名称\(^\text{[47]}\)。
\begin{figure}[ht]
\centering
\includegraphics[width=6cm]{./figures/997dfb34e98d04d5.png}
\caption{约翰·埃勒特·波德，这位提出“天王星”名称的天文学家} \label{fig_Uranus_4}
\end{figure}
在 1782 年 3 月的一篇论文中，约翰·埃勒特·波德提出了“Uranus”这一名称，即希腊天空之神 Ouranos 的拉丁化形式\(^\text{[48]}\)。波德认为，行星名称应遵循神话体系，以免与其他行星显得不同；而 Uranus 作为第一代泰坦之父，是一个合适的名字\(^\text{[48]}\)。他还指出，这个名字十分优雅，因为正如土星（Saturn）是木星之父，新行星也应以土星之父命名\(^\text{[42][48][49][50]}\)。然而，他似乎并不知道 Uranus 只是该神名的拉丁化形式，其真正的罗马对应神是 Caelus。1789 年，波德在皇家科学院的同事马丁·克拉普罗特为其新发现的元素命名为“铀”（uranium），以支持波德的命名选择\(^\text{[51]}\)。最终，波德的建议成为最广泛使用的名称，并在 1850 年成为通用称呼，当时最后坚持使用 Georgium Sidus 的英国航海年鉴局也改用 Uranus\(^\text{[49]}\)。

天王星有两个天文学符号。其一为最早提出者：⛢\(^\text{[h]}\)，由约翰·戈特弗里德·科勒于 1782 年在波德的请求下提出\(^\text{[52]}\)。科勒建议使用铂金的符号，而铂金仅在 30 年前才被科学描述；由于没有传统的炼金术符号，他提出使用 ⛢ 或 ⛢，即由行星–金属符号 ☉（金）与 ♂（铁）组合而来，因为铂（金之“白金”）常与铁共生。波德认为符号直立的版本 ⛢ 更符合其他行星的符号体系，同时依然保持区别性\(^\text{[52]}\)。在现代天文学中，该符号在（罕见的）使用符号的情况下最为常见\(^\text{[53][54]}\)。第二个符号 ♅\(^\text{[i]}\) 是拉朗德于 1784 年提出的；在写给赫歇尔的信中，拉朗德将其描述为“un globe surmonté par la première lettre de votre nom”（“一个上方加上您姓氏首字母的圆球”）\(^\text{[46]}\)。这一符号在占星学中几乎是通用的。

在英语流行文化中，由于 Uranus 的常见发音类似于“your anus”，由此衍生出大量幽默内容\(^\text{[55]}\)。

在其他语言中，天王星有多种名称。在中文（天王星）、日文（天王星 Tennōsei）、韩文（천왕성 Cheonwangseong）与越南语（sao Thiên Vương）中，其名称都直译为“天之王星”\(^\text{[56][57][58][59]}\)。在泰语中，其正式名称为 Dao Yurenat（ดาวยูเรนัส），与英语一致；其另一名称 Dao Maruettayu（ดาวมฤตยู，“Mṛtyu 星”）源于梵语中表示“死亡”的词 mrtyu（मृत्यु）。在蒙古语中，其名称是 Tengeriin Van（Тэнгэрийн ван），意为“天空之王”，反映了其同名神祇作为天界统治者的地位。在夏威夷语中，其名称为 Heleʻekela，即“Herschel”名称的夏威夷语音译\(^\text{[60]}\)。
\subsection{形成}  
一般认为，冰巨行星与气巨行星之间的差异源于它们的形成历史\(^\text{[61][62][63]}\)。太阳系被假设形成于一个被称为原太阳星云的旋转气尘盘中。星云中的大部分气体——主要由氢和氦构成——形成了太阳，而尘埃颗粒聚集在一起形成最初的原行星。随着这些行星体逐渐增长，其中一些最终吸积了足够的物质，使得它们的引力能够束缚住星云中残余的气体\(^\text{[61][62][64]}\)。它们吸附的气体越多，体积就越大；体积越大，又能吸附更多气体，直到达到一个临界点，其大小开始呈指数式增加\(^\text{[65]}\)。冰巨行星仅吸附了数个地球质量的星云气体，因而从未达到这一临界点\(^\text{[61][62][66]}\)。近期关于行星迁移的模拟表明，两颗冰巨行星形成时的位置比现在更靠近太阳，并在形成后向外迁移（即 Nice 模型）\(^\text{[61]}\)。
\subsection{轨道与自转} 
天王星绕太阳一周需要 84 年。自 1781 年被发现以来\(^\text{[67]}\)，依据其相对于背景恒星的位置计算，该行星已经两次回到其最初被发现的位置——在金牛座双星 ζ Tauri 东北方向——分别是在 1865 年 3 月与 1949 年 3 月，并将在 2033 年 4 月再次回到该位置\(^\text{[68]}\)。
\begin{figure}[ht]
\centering
\includegraphics[width=8cm]{./figures/0658c53f45263d76.png}
\caption{在相隔 84 年的两个各为 2 年区间中，天王星的赤经变化} \label{fig_Uranus_5}
\end{figure}
由于其周期非常接近 84 年，因此它在恒星背景中的视位置与整整 84 年前几乎相同（见图）。天王星的逆行运动使其每年都会回到大致与前一年最偏东的位置相近的地方。

它与太阳的平均距离约为 20 AU（30 亿千米；20 亿英里）。其近日点与远日点之间的差值为 1.8 AU，大于任何一颗行星，但仍小于矮行星冥王星的对应差值\(^\text{[69]}\)。太阳光强度与距离的平方成反比——在天王星处（距离约为地球至太阳距离的 20 倍），光照强度约为地球的 1/400\(^\text{[70]}\)。

天王星的轨道要素最早由皮埃尔–西蒙·拉普拉斯于 1783 年计算\(^\text{[71]}\)。随着时间推移，预测轨道与观测轨道之间开始出现偏差，1841 年，约翰·考奇·亚当斯首次提出这些差异可能来自某颗未被观测到的行星的引力扰动。1845 年，乌尔班·勒威耶展开了对天王星轨道的独立研究。1846 年 9 月 23 日，约翰·戈特弗里德·加勒在几乎等同于勒威耶预测的位置发现了一颗新行星，即后来命名的海王星\(^\text{[72]}\)。

天王星内部的自转周期为 17 小时 14 分钟 52 秒\(^\text{[14]}\)，这一数值通过追踪天王星极光的自转运动而得出\(^\text{[73]}\)。与所有巨行星一样，它的高层大气在自转方向上出现强风。在某些纬度（例如大约南纬 60 度），可见大气特征移动得更快，仅需约 14 小时即可完成一次完整自转\(^\text{[74]}\)。
\subsubsection{自转轴倾角}
\begin{figure}[ht]
\centering
\includegraphics[width=8cm]{./figures/a3f72c0a3302a344.png}
\caption{从 1986 年到 2030 年，模拟从地球看到的天王星视图：从 1986 年的南半球夏至，经 2007 年的昼夜平分点，到 2028 年的北半球夏至。} \label{fig_Uranus_6}
\end{figure}
天王星的自转轴大致与太阳系的平面平行，其自转轴倾角可描述为 82.23° 或 97.77°，取决于哪个极被视为北极\({^\text{[j]}}\)。前一种定义遵循国际天文学联合会的规定，即北极是位于太阳系不变平面地球北侧的那一极。在这种定义下，天王星具有逆行自转。或者，若按照另一种惯例，以“旋转方向的右手定则”来定义天体的南北两极，则天王星的自转轴倾角可表示为 97.77°，这将反转哪个极被视为北极、哪个极被视为南极，并使得该行星呈现顺行自转\({^\text{[75]}}\)。这一特性使其季节变化完全不同于其他行星。冥王星与小行星 2 号智神星也具有极端的自转轴倾角。接近至日时，其中一个极持续朝向太阳，另一个极持续背向太阳，只有赤道附近一条狭窄的区域经历快速的昼夜交替，而太阳在地平线上非常低。随着天王星沿轨道移动至另一侧，两极朝向太阳的状态将完全反转。每个极大约经历 42 年的连续白昼，随后是 42 年的连续黑夜\({^\text{[76]}}\)。接近春分或秋分时，太阳照射天王星赤道，使其出现类似其他行星所经历的昼夜交替。

这种自转轴取向的一个结果是：在一个天王星年内，天王星近极地区接收到的太阳能量平均多于赤道地区。然而，天王星的赤道区域却比极地区更热。造成这一现象的内在机制尚不清楚。天王星异常自转轴倾斜的原因也并未确定，但通常的推测是：在太阳系形成期间，一个具有地球大小的原行星与天王星发生碰撞，从而导致了这种偏斜的取向\({^\text{[77]}}\)。达勒姆大学 Jacob Kegerreis 的研究表明，这一倾斜可能是由于一个比地球更大的原行星在 30–40 亿年前撞击天王星所致\({^\text{[78]}}\)。在 1986 年旅行者 2 号探测器飞掠天王星时，其南极几乎正对太阳\({^\text{[79][80]}}\)。

自赫歇尔发现天卫三（Titania）与天卫四（Oberon）以来，天文学家便推测天王星有异常的自转轴，拉普拉斯在 1805 年计算了这些卫星轨道相对于行星赤道平面的倾斜度\({^\text{[81]}}\)。
\begin{figure}[ht]
\centering
\includegraphics[width=8cm]{./figures/733f8b68f8a5af74.png}
\caption{} \label{fig_Uranus_7}
\end{figure}
\subsubsection{从地球的可见性}
\begin{figure}[ht]
\centering
\includegraphics[width=8cm]{./figures/91965dcbef8408da.png}
\caption{“通过业余望远镜观测到的天王星，它出现在 2022 年 11 月月食期间月球掩星之后不久。”} \label{fig_Uranus_8}
\end{figure}
“天王星的平均视星等为 5.68，标准差为 0.17，而其亮度极值为 5.38 和 6.03\(^\text{[20]}\)。这一亮度范围接近肉眼可见的极限。亮度的大部分变化取决于从太阳照亮并从地球观测到的行星纬度\(^\text{[83]}\)。它的视直径介于 3.4 至 3.7 角秒之间，而土星为 16 至 20 角秒，木星为 32 至 45 角秒\(^\text{[84]}\)。在冲日时，天王星在暗夜条件下可被肉眼看到，即使在城市环境中也可以通过双筒望远镜轻松观测到\(^\text{[7]}\)。在口径 15 至 23 厘米的较大业余望远镜上，天王星呈现为一颗浅青色圆盘，并显示明显的边缘变暗。使用 25 厘米或更大口径的望远镜，可观测到云带结构，以及部分较大的卫星，如天卫三（Titania）和天卫四（Oberon）\(^\text{[85]}\)。
\subsubsection{内部结构}
\begin{figure}[ht]
\centering
\includegraphics[width=8cm]{./figures/a3223e5217a60cbf.png}
\caption{地球与天王星的大小比较} \label{fig_Uranus_9}
\end{figure}
天王星的质量约为地球的 14.5 倍，使其成为巨行星中质量最小者。其直径略大于海王星，约为地球的四倍。其由此产生的密度为 1.27 g/cm³，使天王星成为仅次于土星的第二低密度行星\(^\text{[11][12]}\)。这一数值表明它主要由各种“冰”组成，例如水、氨和甲烷\(^\text{[18]}\)。天王星内部冰的总质量并不精确已知，因为不同模型会得出不同数值；其必须介于 9.3 到 13.5 个地球质量之间\(^\text{[18][86]}\)。氢和氦只占总量的一小部分，约为 0.5 至 1.5 个地球质量\(^\text{[18]}\)。剩余的非冰物质（0.5 到 3.7 个地球质量）由岩石物质构成\(^\text{[18]}\)。

天王星的标准结构模型认为它由三层组成：中心为岩质（硅酸盐/铁–镍）核心，中间为冰质地幔，外层为气态的氢/氦包层\(^\text{[18][87]}\)。核心相对较小，质量仅为 0.55 个地球质量，半径不到行星的 20\%；地幔构成其主体，约为 13.4 个地球质量；上层大气相对不显著，质量约 0.5 个地球质量，延伸占天王星半径的最后 20\%\(^\text{[18][87]}\)。天王星核心的密度约为 9 g/cm³，中心压力约为 800 万巴（800 GPa），温度约为 5000 K\(^\text{[86][87]}\)。冰质地幔实际上并非传统意义的固态冰，而是由水、氨及其他挥发物组成的高温高密度流体\(^\text{[18][87]}\)。这种具有高电导率的流体有时被称为“水–氨海洋”\(^\text{[88]}\)。
\begin{figure}[ht]
\centering
\includegraphics[width=8cm]{./figures/9f28c21355c1f0ee.png}
\caption{天王星内部结构示意图，列出了各层的组成。} \label{fig_Uranus_10}
\end{figure}
在天王星深处的极端压力和温度条件下，甲烷分子可能会被分解，其中的碳原子会凝结成钻石晶体，并像冰雹一样穿过地幔降落\({^\text{[89][90]}}\)。这一现象类似于科学家推测可能在木星、土星和海王星上存在的“钻石雨”\({^\text{[91][92]}}\)。劳伦斯利弗莫尔国家实验室的超高压实验表明，在地幔底部可能存在一层金属液态碳的“海洋”，其中也许漂浮着固体“钻石冰山”\({^\text{[93][94][95]}}\)。

天王星和海王星的总体组成与木星和土星不同：冰的含量远多于气体，因此被单独归类为“冰巨行星”。在内部可能存在一层离子态水，其中水分子被分解为氢离子与氧离子的混合物；更深处则可能存在“超离子水”，在这种状态下氧原子形成晶格，而氢离子可以在晶格中自由移动\({^\text{[96]}}\)。

尽管上述结构模型被广泛采用，但并非唯一；其他模型同样能与观测数据相符。例如，如果大量氢和岩质物质混入冰状地幔，那么内部冰的总质量将减少，而相应的岩石和氢的总质量将增加。目前可获得的数据不足以科学地判定哪一种模型是正确的\({^\text{[86]}}\)。天王星内部为流体结构，因此并无固体表面；其气态大气会逐渐过渡到内部的液态层\({^\text{[18]}}\)。为了方便，在大气压等于 1 bar（100 kPa）的位置，约定一个旋转扁球体作为“表面”。该“表面”的赤道半径与极半径分别为 25,559 ± 4 km（15,881.6 ± 2.5 mi）和 24,973 ± 20 km（15,518 ± 12 mi）\({^\text{[11]}}\)。本条目中所有高度皆以此为零点。
\subsubsection{内部热量}
天王星的内部热量显著低于其他巨行星；在天文学中，它被认为具有极低的热通量\({^\text{[24][97]}}\)。为何天王星的内部温度如此之低仍然未被理解。海王星在大小和成分上几乎与天王星相同，但它向外辐射的能量是其从太阳接收能量的 2.61 倍\({^\text{[24]}}\)，而天王星几乎没有额外的热辐射。天王星在远红外（即热辐射）波段释放的总功率是其大气吸收太阳能量的 1.06±0.08 倍\({^\text{[19][16]}}\)。其热通量仅为 0.042±0.047 W/m\(^2\)，甚至低于地球约 0.075 W/m\(^2\) 的内部热通量\({^\text{[16]}}\)。在其对流层顶记录到的最低温度为 49 K（−224.2 °C；−371.5 °F），使天王星成为太阳系中最寒冷的行星\({^\text{[19][16]}}\)。

一种假说认为，引发天王星自转轴倾斜的地球大小撞击体使其核心失去了大量原始热量\({^\text{[98]}}\)。另一种假说认为天王星内部上层存在某种屏障，阻止内部热量向外传输\({^\text{[18]}}\)。例如，具有不同成分的多层结构之间可能发生的对流受抑制，从而减少向上的热传递；\({^\text{[19][16]}}\) 双扩散对流可能是一个限制因素\({^\text{[18]}}\)。

2021 年一项研究通过压缩含橄榄石与铁方镁石等矿物的水来模拟冰巨行星内部环境，结果显示天王星与海王星内部的大量液体中可能溶解了大量镁元素。如果天王星较海王星含有更多镁，则这些镁可能形成隔热层，从而解释该行星异常寒冷的原因\({^\text{[99]}}\)。
\subsection{大气}
尽管天王星内部不存在明确的固体表面，但在遥感可探测范围内，其气态外层的最外部分被称为天王星大气\({^\text{[19]}}\)。遥感探测可以达到 1 bar（100 kPa）层以下约 300 km 的深度，对应压力约 100 bar（10 MPa）及温度约 320 K（47 °C；116 °F）\({^\text{[100]}}\)。稀薄的热层从名义上的 1 bar 表面向外延伸超过两个行星半径\(^\text{[101]}\)。天王星大气通常分为三层：对流层（−300 至 50 km；压力 100 至 0.1 bar）、平流层（50 至 4,000 km；压力 0.1 至 \(10^{-10}\) bar）以及从约 4,000 km 延伸至约 50,000 km 的热层\(^\text{[19]}\)。天王星不存在中气层（mesosphere）。
\subsubsection{成分}
\begin{figure}[ht]
\centering
\includegraphics[width=10cm]{./figures/a69c6ecdefb3fb61.png}
\caption{天王星大气层的组成与结构示意图，以及其压力分布曲线} \label{fig_Uranus_11}
\end{figure}
天王星大气的组成不同于其整体成分，主要由分子氢和氦构成\(^{\text{[19]}}\)。在上对流层中，氦的物质的量分数（即每个气体分子中氦原子的数量）为 \(0.15 \pm 0.03\)\(^{\text{[23]}}\)，对应的质量分数为 \(0.26 \pm 0.05\)\(^{\text{[19][16]}}\)。该数值接近原始太阳氦质量分数 \(0.275 \pm 0.01\)\(^{\text{[102]}}\)，表明氦并未像在气体巨行星中那样沉降至中心\(^{\text{[19]}}\)。天王星大气中第三丰富的成分是甲烷（CH\(_4\)）\(^{\text{[19]}}\)。甲烷在可见光和近红外（IR）区具有显著吸收带，使天王星呈现海蓝色或青色\(^{\text{[19]}}\)。在压力约为 1.3 bar（130 kPa）的甲烷云层以下，甲烷分子在大气中的物质的量分数约为 2.3\%，这一碳丰度约为太阳的 20 至 30 倍\(^{\text{[19][22][103]}}\)。

由于上层大气温度极低，其混合比\(^{\text{[k]}}\)大大降低，使饱和度降低并导致过量甲烷冻结沉降\(^{\text{[104]}}\)。深层大气中较不挥发的化合物（如氨、水和硫化氢）的丰度知之甚少，它们的含量可能也高于太阳丰度\(^{\text{[19][105]}}\)。除了甲烷外，在天王星平流层中还发现了微量多种碳氢化合物，被认为是甲烷在太阳紫外（UV）辐射诱导的光解作用下产生的\(^{\text{[106]}}\)。这些化合物包括乙烷（C\(_2\)H\(_6\)）、乙炔（C\(_2\)H\(_2\)）、甲基乙炔（CH\(_3\)C\(_2\)H）以及二乙炔（C\(_2\)HC\(_2\)H）\(^{\text{[104][107][108]}}\)。光谱分析还揭示了上层大气中存在水蒸气、一氧化碳及二氧化碳的微量成分，这些成分只能来自外源，如尘埃或彗星的输入\({}^{\text{[107][108][109]}}\)。
\subsubsection{对流层}  
对流层是大气中最低、最稠密的部分，其特征是温度随高度升高而降低\(^{\text{[19]}}\)。从标称对流层底部（−300 km处，温度约 320 K（47 °C；116 °F））到 50 km 高度，温度下降到 53 K（−220 °C；−364 °F）\(^{\text{[100][103]}}\)。对流层最冷的上边界区域（即对流层顶）的温度随行星纬度而变化，在 49 至 57 K（−224 至 −216 °C；−371 至 −357 °F）之间\(^{\text{[19][97]}}\)。对流层顶区域负责天王星大部分远红外热辐射，从而决定了其有效温度为 \(59.1 \pm 0.3\) K（−214.1 ± 0.3 °C；−353.3 ± 0.5 °F）\(^{\text{[97][16]}}\)。

对流层被认为具有高度复杂的云结构：理论推测水云位于 50 至 100 bar（5 至 10 MPa）的压力区间，硫化氢铵云位于 20 至 40 bar（2 至 4 MPa），氨或硫化氢云位于 3 至 10 bar（0.3 至 1 MPa），而薄的甲烷云则被直接探测到位于 1 至 2 bar（0.1 至 0.2 MPa）\(^{\text{[19][22][100][110]}}\)。对流层是大气中动态最活跃的部分，表现出强风、明亮云层以及季节变化\(^{\text{[24]}}\)。
\subsubsection{上层大气}
\begin{figure}[ht]
\centering
\includegraphics[width=8cm]{./figures/b45d6aa671a94565.png}
\caption{天王星的高层大气，由哈勃空间望远镜在“外行星大气传承（OPAL）”观测项目中拍摄\(^\text{[111]}\)。} \label{fig_Uranus_12}
\end{figure}
天王星大气的中层是平流层，在这一层中，温度通常随高度升高而升高：从对流层顶处的 53,K（−220,°C；−364,°F）上升到热层底部的 800–850,K（527–577,°C；980–1,070,°F）。$^{\text{[101]}}$ 平流层的加热是由甲烷和其他碳氢化合物吸收太阳紫外与红外辐射所造成的，$^{\text{[112]}}$ 而这些化合物在该大气层中是由甲烷光解作用形成的。$^{\text{[106]}}$ 热量也会从炽热的热层传导下来。$^{\text{[112]}}$ 这些碳氢化合物占据了一个相对狭窄的高度区间，约在 100 至 300,km，对应的气压范围为 1,000 至 10,Pa，温度范围为 75 至 170,K（−198 至 −103,°C；−325 至 −154,°F）。$^{\text{[104][107]}}$ 最丰富的碳氢化合物是甲烷、乙炔与乙烷，其相对于氢气的混合比约为 $10^{-7}$。在这些高度，一氧化碳的混合比也类似。$^{\text{[104][107][109]}}$ 更重的碳氢化合物和二氧化碳的混合比低三个数量级。$^{\text{[107]}}$ 水的丰度约为 $7\times10^{-9}$。$^{\text{[108]}}$ 乙烷和乙炔倾向于在较冷的平流层下部和对流层顶（低于 10,mBar 的区域）凝结，形成薄雾层，$^{\text{[106]}}$ 而这些薄雾可能部分解释了天王星外观缺乏明显特征的原因。薄雾层以上的天王星平流层中，碳氢化合物的浓度明显低于其他巨行星平流层中的浓度。$^{\text{[104][113]}}$
\begin{figure}[ht]
\centering
\includegraphics[width=8cm]{./figures/9588969d08747caa.png}
\caption{天王星——北极——气旋（VLA；2021 年 10 月）} \label{fig_Uranus_13}
\end{figure}
天王星大气的最外层是热层与电晕，其温度在约 800,K（527,°C）至 850,K（577,°C）之间保持均一。$^{\text{[19][113]}}$ 维持如此高温所需的热源尚不清楚，因为无论是太阳紫外辐射还是极光活动都无法提供足够的能量以维持这些温度。平流层中 0.1,mBar 气压层以上缺乏碳氢化合物，使其冷却效率较弱，这也可能有所贡献。$^{\text{[101][113]}}$ 除了分子氢之外，热层—电晕还包含大量自由氢原子。它们的低质量与高温可以解释为何电晕向外延伸至 50,000,km（31,000,mi），即距离天王星表面约两个行星半径的距离。$^{\text{[101][113]}}$

这种高度延展的电晕是天王星独有的特征。$^{\text{[113]}}$ 其影响包括对绕天王星运行的小颗粒产生阻力，从而导致天王星环中的尘埃普遍减少。$^{\text{[101]}}$ 天王星的热层与平流层上部共同构成了天王星的电离层。$^{\text{[103]}}$ 观测显示，电离层分布于 2,000 至 10,000,km（1,200 至 6,200,mi）的高度范围内。$^{\text{[103]}}$ 天王星电离层的密度高于土星或海王星的电离层，这可能源自平流层中碳氢化合物浓度较低。$^{\text{[113][114]}}$ 电离层主要由太阳紫外辐射维持，其密度随太阳活动而变化。$^{\text{[115]}}$ 与木星和土星相比，天王星的极光活动微不足道。$^{\text{[113][116]}}$
\subsection{气候}
主条目：天王星的气候
在紫外与可见光波段，天王星的大气相比其他巨行星（甚至与其极为相似的海王星相比）都显得较为平淡。$^{\text{[24]}}$ 当旅行者 2 号飞掠天王星（1986 年）时，在整颗行星上仅观测到 10 个云系特征。$^{\text{[117][118]}}$ 对此特征稀少的一种解释是天王星内部热量显著低于其他巨行星，使其成为太阳系中最冷的行星。$^{\text{[19][16]}}$
\subsubsection{带状结构、风与云层}
\begin{figure}[ht]
\centering
\includegraphics[width=8cm]{./figures/d19b68635fe857ad.png}
\caption{“旅行者2号拍摄的天王星动态大气延时影像”} \label{fig_Uranus_14}
\end{figure}
1986 年，旅行者 2 号发现天王星可见的南半球可以分为两个区域：明亮的极冠与黑暗的赤道带。$^{\text{[117]}}$ 它们的边界位于大约 −45° 纬度处。在 −45° 至 −50° 的纬度范围内有一条狭窄的带状结构，是其可见表面上最明亮的大型特征。$^{\text{[117][119]}}$ 它被称为南部“项圈”。极冠与项圈被认为是位于 1.3 至 2,bar 气压范围内的致密甲烷云区域。$^{\text{[120]}}$ 除了大尺度的带状结构外，旅行者 2 号还观测到了十个小型明亮云团，大多位于项圈以北数度的位置。$^{\text{[117]}}$ 在其他方面，1986 年的天王星看起来像是一颗动力学上死寂的行星。

旅行者 2 号到达天王星时正值其南半球盛夏，因此无法观测北半球。进入 21 世纪初，当北极区域开始进入视野时，哈勃太空望远镜（HST）与凯克望远镜最初在北半球既未观测到项圈，也未观测到极冠。$^{\text{[119]}}$ 因此天王星呈现出一种不对称性：在南极附近明亮，而在南部项圈以北的区域则均匀黑暗。$^{\text{[119]}}$ 2007 年，当天王星通过其春分点时，南部项圈几乎消失，而在约 45° 纬度附近出现了一条暗淡的北部项圈。$^{\text{[121]}}$ 2023 年，一个利用甚大天线阵（VLA）的研究团队在 80° 纬度观测到一条黑暗的项圈，并在北极发现一个明亮斑点，指示存在一个极涡。$^{\text{[122]}}$
\begin{figure}[ht]
\centering
\includegraphics[width=8cm]{./figures/ab9d610ba2b86b0e.png}
\caption{天王星上观测到的第一个暗斑。图像由哈勃太空望远镜 ACS 于 2006 年获取。} \label{fig_Uranus_15}
\end{figure}
在 1990 年代，观测到的明亮云系数量显著增加，部分原因是新的高分辨率成像技术变得可用。$^{\text{[24]}}$ 随着北半球开始进入视野，大多数明亮云系都被发现位于北半球。$^{\text{[24]}}$ 一个早期的解释——即明亮云系在天王星较暗的北半球更容易被识别，而在南半球，明亮的项圈会遮蔽它们——已被证明是不正确的。$^{\text{[123][124]}}$ 尽管如此，两半球的云系之间确实存在差异。北半球的云更小、更锐利、也更明亮。$^{\text{[124]}}$ 它们似乎位于更高的高度。$^{\text{[124]}}$ 云系的寿命跨度跨越数个数量级：一些小型云团仅存在数小时；而至少有一团南半球云可能自旅行者 2 号飞掠以来一直存在。$^{\text{[24][118]}}$ 最近的观测还发现，天王星上的云系特征与海王星上的有许多共同之处。$^{\text{[24]}}$ 例如，海王星上常见的暗斑在 2006 年之前从未在天王星上被观测到，当时首次拍摄到的此类结构被命名为“天王星暗斑”。$^{\text{[125]}}$ 有一种推测认为，天王星正在其春分季节中变得越来越类似海王星。$^{\text{[126]}}$

对大量云系特征的追踪使得能够确定天王星上对流层上部的纬向风。$^{\text{[24]}}$ 在赤道附近，风向为逆行，即与行星自转方向相反，其风速介于 −360 至 −180,km/h（−220 至 −110,mph）之间。$^{\text{[24][119]}}$ 随着远离赤道，风速增大，在约 ±20° 纬度附近（对流层温度最低处）接近零。$^{\text{[24][97]}}$ 更接近两极时，风向转为顺行，与天王星的自转方向一致；风速继续增加，在 ±60° 纬度达到最大值，然后在极地下降至零。$^{\text{[24]}}$ 在 −40° 纬度，风速范围为 540 至 720,km/h（340 至 450,mph）。由于南部项圈会遮蔽其以南的所有云系，因此项圈与南极间的风速无法测量。$^{\text{[24]}}$ 相比之下，在北半球，观测到的最大风速可达 860,km/h（540,mph），出现在约 +50° 纬度附近。$^{\text{[24][119][127]}}$

1986 年，旅行者 2 号的行星射电天文（PRA）实验观测到了 140 次闪电，或称天王星静电放电，其频率在 0.9–40,MHz 范围内。$^{\text{[128][129]}}$ 在 24 小时内，这些 UEDs（静电放电）在距离天王星 600,000,km 的位置被探测到，其中大多数不可见。$^{\text{[128]}}$ 然而，微物理模型表明，天王星的闪电发生在对流层深处的水云对流风暴中。$^{\text{[128][130]}}$ 若情况如此，则闪电将不会可见，因为其上方存在厚厚的云层。$^{\text{[129]}}$ 天王星闪电的功率约为 10^8,W，释放 $1\times10^7-2\times10^7,J$ 的能量，平均持续约 120,ms。有可能天王星闪电的功率会随季节显著变化，这可能由云层中对流速率的变化引起。$^{\text{[129]}}$ 天王星闪电比地球闪电强得多，其强度与木星闪电相当。$^{\text{[129]}}$ 在冰巨星飞掠期间，“旅行者 2 号”探测到的天王星闪电比海王星更清晰，这可能是由于天王星较低的重力以及可能更温暖的深层大气所致。$^{\text{[130]}}$
\subsubsection{季节变化}
\begin{figure}[ht]
\centering
\includegraphics[width=6cm]{./figures/da737c10e7b53f23.png}
\caption{2005 年的天王星。可见其光环、南部项圈以及北半球的一处明亮云团（哈勃太空望远镜 ACS 图像）。} \label{fig_Uranus_16}
\end{figure}
从 2004 年 3 月到 5 月的一段短暂时期内，天王星大气中出现了大片云层，使其呈现出类似海王星的外观。$^{\text{[131][124][132]}}$ 这些观测包括创纪录的 820,km/h（510,mph）风速，以及一个被称为“七月四日烟火”的持久雷暴。$^{\text{[118]}}$ 2006 年 8 月 23 日，科罗拉多州博尔德的空间科学研究所与威斯康星大学的研究人员在天王星表面观测到一个暗斑，使科学家对天王星的大气活动有了更多了解。$^{\text{[125]}}$ 这一突发活动上升的原因尚不完全清楚，但似乎与天王星极端的自转轴倾角有关，这导致其天气出现极端季节性变化。$^{\text{[133][126]}}$ 确定这种季节变化的性质很困难，因为有关天王星大气的优质数据存在时间不足 84 年，也就是一个完整的天王星年。自 1950 年代以来跨越半个天王星年的光度测量显示在两个光谱波段上存在规律性亮度变化，其最大值出现在至日附近，最小值出现在分点附近。$^{\text{[134]}}$ 自 1960 年代开始的对深层对流层微波辐射的测量也记录到了类似的周期性变化，其最大值同样出现在至日。$^{\text{[135]}}$ 20 世纪 70 年代开始的平流层温度测量也显示，在 1986 年至日前后出现了温度峰值。$^{\text{[112]}}$ 这一变化的主要原因被认为是观测几何视角的改变。$^{\text{[123]}}$

有迹象表明天王星上也确实发生了物理性的季节变化。尽管天王星已知拥有明亮的南极区域，但其北极却相当昏暗，这与先前所描述的季节变化模型并不一致。$^{\text{[126]}}$ 在 1944 年前一次北半球至日时，天王星表现出更高的亮度，说明北极并非一直如此昏暗。$^{\text{[134]}}$ 这些信息意味着可见极区会在至日前一段时间变亮，并在分点后变暗。$^{\text{[126]}}$ 对可见光与微波数据的详细分析显示：这些周期性亮度变化并未在至日前后呈现完全对称性，这也表明经向反照率结构正在发生变化。$^{\text{[126]}}$

在 1990 年代，随着天王星远离其至日，哈勃望远镜与地基望远镜的观测显示南极极冠区域明显变暗（明亮的南部项圈除外），$^{\text{[120]}}$ 而北半球则表现出活动增强，$^{\text{[118]}}$ 例如形成云系与更强风速，使人们预期其将很快变亮。$^{\text{[124]}}$ 这一变化确实发生在 2007 年天王星经过分点时：一个暗淡的北部极项圈出现，而南部项圈几乎变得不可见，尽管纬向风速分布仍略微不对称，北半球风速较南半球稍慢。$^{\text{[121]}}$

这些物理变化的机制仍不明确。$^{\text{[126]}}$ 在夏至与冬至附近，天王星的半球会交替完全暴露在太阳照射下，或完全面对深空。被阳光照亮的半球变亮，被认为是由于对流层中局部甲烷云与薄雾层变得更为厚实。$^{\text{[120]}}$ 位于 −45° 纬度的明亮项圈也与甲烷云有关。$^{\text{[120]}}$ 南极区域的其他变化则可由较低云层结构的变化解释。$^{\text{[120]}}$ 天王星微波辐射的变化很可能由深层对流层环流的变化引起，因为厚实的极地区云与薄雾层可能抑制对流。$^{\text{[136]}}$ 随着天王星的春分与秋分到来，大气动力学正在发生改变，对流可以再次出现。$^{\text{[118][136]}}$
\subsection{磁层}
\begin{figure}[ht]
\centering
\includegraphics[width=10cm]{./figures/17771a29d74a5a06.png}
\caption{天王星的磁场（动画；2020 年 3 月 25 日）} \label{fig_Uranus_17}
\end{figure}
在旅行者 2 号抵达之前，从未对天王星磁层进行过测量，因此其性质一直是个谜。1986 年之前，科学家原本预计天王星的磁场应当与太阳风方向一致，因为在这种情况下磁场会与位于黄道面中的天王星自转极对齐。$^{\text{[137]}}$

旅行者号的观测显示，天王星的磁场极为特殊，不仅因为其磁场并非起源于行星的几何中心，而且其磁轴还相对于自转轴倾斜了 59°。$^{\text{[137][138]}}$ 实际上，磁偶极子从天王星中心向南旋转极方向偏移了多达行星半径的三分之一。$^{\text{[137]}}$ 这种异常的几何结构导致其磁层呈现高度不对称：在南半球表面，磁场强度最低可降至 0.1,gauss（10,μT），而在北半球表面则最高可达 1.1,gauss（110,μT）。$^{\text{[137]}}$ 天王星表面的平均磁场强度为 0.23,gauss（23,μT）。$^{\text{[137]}}$
\begin{figure}[ht]
\centering
\includegraphics[width=8cm]{./figures/77da8bf7c1d01605.png}
\caption{一幅展示天王星非对称磁层的示意图} \label{fig_Uranus_18}
\end{figure}
对旅行者 2 号数据的 2017 年研究表明，这种不对称性导致天王星的磁层在每个天王星日都会与太阳风相连，使得行星暴露在来自太阳的粒子之下。$^{\text{[139]}}$ 相比之下，地球的磁场在两极处强度大致相同，其“磁赤道”也与地理赤道大致平行。$^{\text{[138]}}$ 天王星的磁偶极矩是地球的 50 倍。$^{\text{[137][138]}}$ 海王星同样具有偏移且倾斜的磁场，暗示这可能是冰巨行星的共同特征。$^{\text{[138]}}$ 一种假说认为，与类地行星和气体巨行星的磁场在其核心内产生不同，冰巨行星的磁场可能由相对浅层的流体运动产生，例如在水–氨海洋中。$^{\text{[88][140]}}$ 另一个可能的解释是天王星内部存在液态钻石的海洋，它们会影响磁场的排列。$^{\text{[94]}}$

然而，目前尚不清楚天王星磁场所呈现的这种不对称性是否代表其磁层的典型状态，或仅仅是由于在异常的太空天气条件下被观测到的巧合。对 2024 年旅行者数据的后续分析显示，飞掠期间观测到的磁层强烈不对称形状代表一种异常状态，因为当时测得的太阳风密度异常偏高，这可能压缩了天王星的磁层。与太阳风事件的相互作用还可能解释一个显著的悖论：尽管测得的磁层等离子体密度普遍偏低，但却存在强烈的电子辐射带。这类条件估计只在不到 5\% 的时间里发生。$^{\text{[141][142]}}$

尽管磁层的排列方式颇为奇特，但在其他方面天王星的磁层与其他行星类似：其前方约 23 个天王星半径处有弓形激波，其磁层顶位于 18 个天王星半径处，具有完全发展的磁尾和辐射带。$^{\text{[137][138][143]}}$ 总体而言，天王星磁层的结构不同于木星，更类似于土星。$^{\text{[137][138]}}$ 天王星的磁尾在其身后延伸数百万公里，由于其横躺的自转方式而被扭曲成一条长长的螺旋形。$^{\text{[137][144]}}$
\begin{figure}[ht]
\centering
\includegraphics[width= 8cm]{./figures/2926db344ae9b351.png}
\caption{由安装在哈勃空间望远镜上的空间望远镜成像光谱仪（STIS）拍摄的天王星极光。$^{\text{[145]}}$} \label{fig_Uranus_19}
\end{figure}
天王星的磁层中包含带电粒子：主要是质子和电子，还有少量的 $H(_2^+)$ 离子。$^{\text{[138][143]}}$ 这些粒子中许多可能来源于热层。$^{\text{[143]}}$ 离子和电子的能量可分别高达 4 与 1.2,MeV。$^{\text{[143]}}$ 内磁层中低能（低于 1,keV）离子的密度约为 $2,cm(^{-3})$。$^{\text{[146]}}$ 粒子群体受天王星卫星的强烈影响，这些卫星在穿过磁层时会扫出明显的空隙。$^{\text{[143]}}$ 粒子通量足够高，可在极短的天文时间尺度（约 10 万年）内使其表面变暗或发生空间风化。$^{\text{[143]}}$ 这可能是天王星卫星与光环整体呈现暗色的原因。$^{\text{[147]}}$

天王星拥有相对发达的极光，其表现为围绕南北磁极的明亮弧形结构。$^{\text{[113]}}$ 不同于木星，天王星的极光似乎对维持行星热层的能量平衡并无显著贡献。$^{\text{[116]}}$ 天王星极光（更准确说，是其三氢阳离子 $H(_3^+)$ 的红外谱线辐射）已在 2023 年底之前被深入研究。$^{\text{[148]}}$

2020 年 3 月，美国国家航空航天局（NASA）的天文学家在重新分析旅行者 2 号在 1986 年飞掠天王星时记录的旧数据后，报告探测到一个大型大气磁泡，也称为“等离子团”（plasmoid），被天王星释放到外太空中。$^{\text{[149][150]}}$
\subsection{卫星}
\begin{figure}[ht]
\centering
\includegraphics[width=8cm]{./figures/bb76da43ae798b6b.png}
\caption{按与天王星距离由近及远（从左到右）排列的主要卫星，并按其真实相对大小与反照率展示。从左到右依次为米兰达、艾丽尔、安布里尔、泰坦妮娅和奥伯龙。（旅行者 2 号拍摄照片的拼接图）} \label{fig_Uranus_20}
\end{figure}
天王星已知有 29 颗天然卫星。$^{\text{[151][152]}}$ 这些卫星的名称选自威廉·莎士比亚与亚历山大·蒲柏作品中的人物。$^{\text{[87][153]}}$ 五颗主要卫星分别是米兰达、艾丽尔、安布里尔、泰坦妮娅和奥伯龙。$^{\text{[87]}}$ 在所有巨行星中，天王星的卫星系统质量最小；五颗主要卫星的总质量还不到海王星最大卫星——海卫一（特里同）单独一颗的一半。$^{\text{[12]}}$ 天王星最大的卫星泰坦妮娅，其半径仅为 788.9,km（490.2,mi），不到月球的一半，但略大于土星第二大卫星瑞亚，因此泰坦妮娅是太阳系中第八大卫星。天王星卫星的反照率相对较低，从安布里尔的 0.20 到艾丽尔的 0.35（在绿光波段）。$^{\text{[117]}}$ 它们是由大约 50\% 的冰与 50\% 的岩石组成的冰–岩混合体。冰可能包含氨与二氧化碳。$^{\text{[147][154]}}$

在天王星卫星之中，艾丽尔的表面似乎最年轻，碰撞坑最少，而安布里尔最古老。$^{\text{[117][147]}}$ 米兰达具有深达 20,km（12,mi）的断层峡谷、阶状层结构，以及表面年龄与特征呈混乱变化的区域。$^{\text{[117]}}$ 米兰达过去的地质活动被认为由潮汐加热驱动，当时其轨道偏心率比现在更大，可能源于其曾与安布里尔经历过 3:1 的轨道共振。$^{\text{[155]}}$ 与上涌的地幔团块（diapirs）相关的伸展作用很可能是米兰达那些类似“赛车跑道”的冠状结构（coronae）的起源。$^{\text{[156][157]}}$ 认为艾丽尔曾经与泰坦妮娅处于 4:1 的轨道共振中。$^{\text{[158]}}$

天王星至少拥有一个占据太阳–天王星 L3 拉格朗日点的马蹄轨道天体——即 83982 Crantor，该点是一处在轨道上 180° 位置、引力不稳定的区域。$^{\text{[159][160]}}$ Crantor 在天王星的同轨区内沿着复杂的、暂时性的马蹄形轨道运行。2010 EU65 也是一个有潜力的天王星马蹄共振体候选者。$^{\text{[160]}}$
\begin{figure}[ht]
\centering
\includegraphics[width=6cm]{./figures/920c89c27309c1b5.png}
\caption{詹姆斯·韦布太空望远镜 NIRCam 拍摄的天王星及其五颗主要卫星和九颗内侧卫星} \label{fig_Uranus_21}
\end{figure}
\subsection{光环}
\begin{figure}[ht]
\centering
\includegraphics[width=8cm]{./figures/6386ea466f33d972.png}
\caption{由詹姆斯·韦布空间望远镜的近红外相机拍摄的天王星光环、内侧卫星与大气层。} \label{fig_Uranus_22}
\end{figure}
天王星的光环由极其暗淡的粒子组成，其尺寸范围从微米级到不足一米。$^{\text{[117]}}$ 目前已知有十三条独立的光环，其中最明亮的是 ε 光环。除两条光环外，天王星的所有光环都极其狭窄——通常只有数公里宽。这些光环可能相当年轻；动力学分析表明它们并非与天王星同时形成。光环物质可能曾属于被高速撞击打碎的一颗（或多颗）卫星。由撞击产生的众多碎片中，只有少量粒子在具有稳定性的轨道区域存留下来，形成了现今的光环。$^{\text{[147][161]}}$

威廉·赫歇尔在 1789 年描述过天王星周围可能存在的光环。这一观测一般被认为存疑，因为光环非常暗淡，并且接下来两个世纪中没有其他天文学家再度观测到它们。然而，赫歇尔对 ε 光环的尺寸、相对于地球的倾角、红色外观以及天王星绕太阳运行时其外观变化的描述却非常准确。$^{\text{[162][163]}}$ 光环系统最终在 1977 年 3 月 10 日由詹姆斯·L·埃利奥特、爱德华·W·邓纳姆和杰西卡·明克利用柯伊伯空中天文台最终确认。他们的发现带有偶然性：原本计划利用恒星 SAO 158687（也称 HD 128598）被天王星掩星的事件研究其大气结构。分析观测数据后，他们发现恒星在消失于天王星之后，在其前后都曾短暂消失五次。他们因此得出天王星周围一定存在光环系统的结论。$^{\text{[164]}}$ 后来，他们又探测到了另外四条光环。$^{\text{[164]}}$ 当旅行者 2 号在 1986 年飞掠天王星时，光环首次被直接成像。$^{\text{[117]}}$ 旅行者 2 号还发现了两条额外的暗淡光环，使光环总数达到十一条。$^{\text{[117]}}$

2005 年 12 月，哈勃太空望远镜探测到一对此前未知的光环。其中最大的一条位于天王星先前已知光环距离的两倍之外。这些新光环由于距离行星极远，被称为“外光环系统”。哈勃还发现了两颗小卫星，其中之一——马布（Mab）——与最外侧的新光环共轨。新光环的加入使天王星光环总数达到十三条。$^{\text{[165]}}$ 2006 年 4 月，来自凯克天文台的新光环图像揭示了外光环的颜色：最外侧那一条呈深蓝色，而另一条呈红色。$^{\text{[166][167]}}$ 关于外侧光环呈蓝色的一种假说是，它由来自马布表面的微小水冰粒子组成，这些粒子足够细小，能够散射蓝光。$^{\text{[166][168]}}$ 相比之下，天王星的内侧光环呈灰色。$^{\text{[166]}}$

尽管天王星光环从地球上非常难以直接观测，但数字成像技术的发展使得一些业余天文学家利用红光或红外滤镜成功拍摄到了光环；如果具备合适的成像设备，口径仅 36,cm（14 英寸）的望远镜也有可能检测到光环。$^{\text{[169]}}$
\subsection{探索}
\begin{figure}[ht]
\centering
\includegraphics[width=6cm]{./figures/410941405bf9f85e.png}
\caption{从土星轨道上的卡西尼号飞船观测到的天王星。} \label{fig_Uranus_23}
\end{figure}
1977 年发射的旅行者 2 号在 1986 年 1 月 24 日最接近天王星，距离云顶仅 81,500,km（50,600,mi），随后继续前往海王星。该航天器研究了天王星大气的结构与化学成分，$^{\text{[103]}}$ 包括由于其极端自转轴倾角而产生的独特天气系统。它首次对天王星五颗最大卫星进行了详细调查，并发现了另外 10 颗新卫星。旅行者 2 号检查了系统中已知的九条光环，并发现了两条新的光环。$^{\text{[117][147][170]}}$ 它还研究了天王星磁场、其不规则结构、倾角，以及由天王星横躺式自转导致的独特螺旋状磁尾。$^{\text{[137]}}$

自那以后，再无其他航天器飞掠天王星，尽管已经提出过许多重访天王星系统的任务方案。在 2009 年的任务延长期规划阶段中，科学家评估了将“卡西尼”号从土星送往天王星的可能性，但最终选择让其进入土星大气销毁，$^{\text{[171]}}$ 因为若从土星出发，需要约 20 年才能抵达天王星系统。$^{\text{[171]}}$ 一个进入天王星大气的探测器可能会采用“金星先驱多探测器”（Pioneer Venus Multiprobe）的技术传承，并下降到 1–5 个大气压区域。$^{\text{[172]}}$ 2011 年发布的《2013–2022 行星科学十年调查》推荐了“天王星轨道器与探测器”任务；该方案设想于 2020–2023 年发射，并进行长达 13 年的巡航前往天王星。$^{\text{[172]}}$ 这一优先级在 2022 年得到再次确认，当时天王星探测器/轨道器任务被列为最高优先事项，原因是人类对冰巨行星的认识极为有限。$^{\text{[173]}}$ 最新的方案是中国国家航天局（CNSA）的“天问四号”木星轨道器，预计 2029 年发射，其携带的子探测器将在任务中脱离，利用木星重力助推而非入轨，在 2045 年 3 月飞掠天王星，随后前往星际空间。$^{\text{[26]}}$ 中国还提出了潜在的“天问五号”任务方案，可能选择环绕天王星或海王星，但尚未正式立项执行。$^{\text{[26]}}$
\subsection{文化影响}
除了在小说作品中频繁出现之外，天王星还启发了许多艺术创作，包括莉迪亚·西古妮（Lydia Sigourney）1827 年的诗歌《乔治亚行星》（*The Georgian Planet*），以及古斯塔夫·霍尔斯特（Gustav Holst）在 1914–1916 年间创作的管弦乐组曲《行星组曲》（*The Planets*）中的一个乐章。赫歇尔发现天王星也被约翰·济慈（John Keats）在诗作《初读查普曼译荷马》（“On First Looking into Chapman’s Homer”）中提及：“那时我恰似仰望苍穹的守望者，有一颗新行星游入他的视野。”$^{\text{[174]}}$ 天王星的发现还启发了化学元素铀（uranium）的命名，该元素由德国化学家马丁·海因里希·克拉普罗特（Martin Heinrich Klaproth）于 1789 年发现。$^{\text{[175]}}$

在现代占星学中，天王星（符号：Uranus monogram）是水瓶座的守护星；在发现天王星之前，水瓶座的守护星是土星。由于天王星呈青蓝色，并与电有关，因此接近青色的“电蓝色”（electric blue）被视为水瓶座的代表色。$^{\text{[176]}}$
\subsection{另见}
\begin{itemize}
\item 2011 QF99 与 2014 YX49，已知仅有的两颗天王星特洛伊小行星
\item 天王星的殖民
\item 地外钻石（被认为在天王星中十分丰富）
\item 天王星概览
\item 太阳系各行星统计
\item 占星学中的天王星
\item 文学与虚构作品中的天王星”
\end{itemize}
\subsection{备注}\\
a.基于欧文，帕特里克·G·J; 多宾森，杰克; 詹姆斯，阿朱那; 提比，尼古拉斯·A; 西蒙，艾米·A; 弗莱彻，利·N; 罗曼，迈克尔·T; 奥顿，格伦·S；黄，迈克尔·H; 托莱多，丹尼尔; 佩雷斯-霍约斯，圣地亚哥; 贝克，朱莉 (2023年12月23日)。“模拟天王星的颜色和大小的季节性周期，并与海王星进行比较”。皇家天文学会月刊。527(4):11521-11538。doi:10.1093/mnras/stad3761。高密度脂蛋白:20.500.11850/657542。ISSN 0035-8711。\\
b.这些是来自VSOP87的平均元素以及导出的量。\\
c.指1巴大气压的水平。\\
d.使用Seidelmann，2007的数据计算。$^{\text{[11]}}$\\
c.基于1 bar大气压水平内的体积。\\
f.He，H的计算2和CH4摩尔分数基于甲烷与氢气的2.3\% 混合比和15/85 He/H2在对流层顶测量的比例。\\
g.因为，在英语世界，后者听起来像 “你的肛门”，以前的发音也省去了尴尬: 作为Pamela同性恋,一位天文学家南伊利诺伊大学爱德华兹维尔,在她的播客上注明，以避免 “被任何小学生取笑……当有疑问时，不要强调任何事情，只要说/ˈjʊərənəs/。然后跑，快。$^{\text{[43]}}$\\
h.Cf.⛢(并非所有字体都支持)\\
i.Cf.♅(并非所有字体都支持)\\
j.匹配页面上的描述轴向倾斜,82.23 ° 而不是97.77 ° 在整个页面中使用。请参见以下解释。\\
k.混合比定义为每分子氢的化合物的分子数。










